% Copyright (c) 2026 Christian Reitwiessner. Apache-2.0.
% Authors: Christian Reitwiessner.
\documentclass[10pt,a4paper]{article}
\usepackage[margin=23mm]{geometry}
\usepackage{amsmath,amssymb,booktabs,hyperref}
\hypersetup{colorlinks=true,urlcolor=blue,pdftitle={A public-message privacy counterexample: STOC submission 1234},pdfauthor={Christian Reitwiessner}}
\setlength{\parindent}{0pt}
\setlength{\parskip}{6pt}
\newcommand{\Prb}{\mathop{\mathrm{Pr}}}
\title{A public-message privacy counterexample\\\large STOC submission 1234: the reverse Tetrahedral protocol}
\author{Christian Reitwiessner}
\date{6 November 2026}
\begin{document}
\maketitle
\begin{abstract}
The reverse protocol in Theorem 5.1 of Donald Beaver's \emph{Correlated Group Interactions: Founding Cryptography on $S_4$ and $S_n$} leaks information about Alice's hand through its public message. Six equiprobable cases suffice to exhibit the failure. This note gives the exact calculation and its correspondence to a checked Lean theorem.
\end{abstract}
\section{Experiment and privacy requirement}
The source is Theorem 5.1, PDF page 5, Figures 9 and 10 of \href{https://github.com/sigact-stoc/stoc2027/releases/download/papers-2/1234.pdf}{the submission manuscript}. Alice chooses an ordered pair $A=(a,b)$ of distinct cards from $\{0,1,2,3\}$ uniformly among the twelve possibilities. She publishes $M$ according to the Alice table. An independent uniform OT selection bit determines Bob's card $C$ via Figure 9. For each $A$, its two OT cases yield the two cards outside $A$, once each. Thus $(a,b,C,D)$ is uniform over all 24 permutations, where $D$ is the discarded card.

In an ideal $(2,1;1)$ deal, Bob receives only $C$. In the real protocol he also sees $M$. A necessary privacy condition is that this public component of his view can be simulated from $C$ alone, preserving its joint distribution with Alice's hand. In particular,
$$
 \Prb[A=a\mid C=c,M=m]=\Prb[A=a\mid C=c]
$$
whenever the conditioning event has positive probability. Failure for $(C,M)$ suffices: Bob's other view components need not be examined.
\section{Six equiprobable cases}
Condition on $C=0$. The following six Figure 9 cases each have conditional probability $1/6$:
\begin{center}
\begin{tabular}{cccc}
\toprule Alice's ordered hand $A$ & Bob $C$ & Discard $D$ & Public $M$\\\midrule
$(1,2)$ & 0 & 3 & $Y$\\
$(1,3)$ & 0 & 2 & $X$\\
$(2,1)$ & 0 & 3 & $X$\\
$(2,3)$ & 0 & 1 & $Z$\\
$(3,1)$ & 0 & 2 & $Z$\\
$(3,2)$ & 0 & 1 & $Y$\\\bottomrule
\end{tabular}
\end{center}
Before the message, Bob gives probability $1/6$ to $A=(1,3)$. Once he sees $M=X$, only $(1,3)$ and $(2,1)$ remain, each with probability $1/2$.
\newpage
\section{Exact leakage calculation}
We obtain the strict inequality
$$
 \Prb[A=(1,3)\mid C=0]=\frac16
 \ne \frac12=\Prb[A=(1,3)\mid C=0,M=X].
$$
The defect persists even if the order of Alice's cards is ignored. With $H=\{a,b\}$:
\begin{center}
\begin{tabular}{ccc}
\toprule Alice's unordered hand $H$ & $\Prb[H\mid C=0]$ & $\Prb[H\mid C=0,M=X]$\\\midrule
$\{1,2\}$ & $1/3$ & $1/2$\\
$\{1,3\}$ & $1/3$ & $1/2$\\
$\{2,3\}$ & $1/3$ & $0$\\\bottomrule
\end{tabular}
\end{center}
In particular, Bob rules out the previously possible hand $\{2,3\}$. A simulator given only $C$ cannot reproduce this extra correlation with Alice's hand.
\section{The checked Lean counterexample}
Let $E$ denote $A=(1,3)$, $F$ denote $C=0$, and $G$ denote $M=X$. Among the 24 shuffles,
$$
 |E\cap F\cap G|=1,\qquad |F|=6,\qquad
 |E\cap F|=1,\qquad |F\cap G|=2.
$$
Conditional independence would require
$$
 |E\cap F\cap G|\,|F|=|E\cap F|\,|F\cap G|.
$$
Instead, the two sides are $6$ and $2$. The declaration
\begin{center}
\texttt{Cslib.Crypto.Tetrahedral.reverse\_message\_leakage}
\end{center}
checks precisely this division-free inequality, using \texttt{by decide} to enumerate the finite cases. Separate checked declarations establish that the enumerator consists of exactly the 24 distinct permutations. The message encoding is $X=0$, $Y=1$, $Z=2$.
\section{Faulty inference and scope}
The proof of Theorem 5.1 argues from $M$ being determined by Alice's hand to privacy. Determinism given Alice's hand does not imply independence from that hand given Bob's card. Checking a joint probability only for a message consistent with a fixed hand does not prove the required conditional independence; incompatible hands are exactly what the observed message excludes.

The counterexample contradicts perfect semi-honest privacy for this reverse protocol. It does not refute its dealt-output distribution, the forward protocol's checked correctness/privacy, or the existence of a different secure reverse construction. It is not a claim that the entire paper is false.
\section*{Sources and reproducibility}
\small
Primary source: manuscript PDF page 5, Figures 9--10 and Theorem 5.1 (linked on page 1). Protocol tables are also Figures 2--3. The \href{https://github.com/cslib-community/complexitylib/blob/4d578733628a41dfc7f550f80725960b87de3c4a/Complexity/Crypto/Multiparty/TetrahedralYield.lean#L106}{immutable Lean proof} is in \texttt{TetrahedralYield.lean}, lines 106--116, in \href{https://github.com/cslib-community/complexitylib/pull/2529}{Complexitylib PR 2529}.

Revision: \texttt{4d578733628a41dfc7f550f80725960b87de3c4a}.

Kernel checking establishes the counting inequality; correspondence to the manuscript additionally uses the explicit table mapping above. AI-assisted technical note; copyright 2026 Christian Reitwiessner, Apache-2.0.
\end{document}
